\documentclass[11pt]{article}
\usepackage[margin=1in]{geometry}
\usepackage{amsmath,amssymb,amsthm,booktabs,tabularx}
\usepackage[T1]{fontenc}
\usepackage{lmodern}
\usepackage[hidelinks]{hyperref}
\setlength{\emergencystretch}{2em}
\newtheorem{lemma}{Lemma}
\newtheorem{theorem}[lemma]{Theorem}
\newtheorem{corollary}[lemma]{Corollary}
\newcommand{\qi}[1]{[#1]_q}
\newcommand{\coef}[2]{[q^{#1}]#2}
\title{Coefficient injections and equality for Sagan--Swanson bottom heaviness\\
\large An alternative proof after Ma and Principia Math}
\author{Haobo Ma\\\small ChronoAI Pte Ltd, Singapore
\and Wenlin Zhang\\\small National University of Singapore\\\small ChronoAI Pte Ltd, Singapore}
\date{November 14, 2026}
\begin{document}
\maketitle
\begin{abstract}
We give an alternative proof, after Shi-Mei Ma (2026) and Principia
Math in Progress (2026, Theorem 1), of Sagan and Swanson's
Conjecture 7.6 (arXiv:2205.14078v1, Section 7.5, p.~40, posed
May 27, 2022), using concentric interval layers for all four unshifted families.
All four unshifted $q$-Stirling families are bottom heavy: each
coefficient below the middle is at least its reflected coefficient.
Concentric interval layers give an elementary proof for every
nonnegative linear combination of products of $q$-integers.
The layers give coefficient injections, an exact formula for the
reflection defect, and a classification of symmetric entries.
We deduce complex-reflection-group consequences and a closed
alternating-sum formula clarifying the symmetry in the source's Conjecture 7.9.
\end{abstract}

\section{The statement and its conventions}
Sagan and Swanson ask whether the coefficients of $S[n,k]$, $c[n,k]$,
$S_B[n,k]$, and $c_B[n,k]$ are bottom heavy \cite[Conjecture 7.6]{SS}.
For a nonzero polynomial $P(q)=\sum_{i=0}^D a_iq^i$ of actual degree $D$,
this means
\begin{equation}\label{eq:bh}
 a_{D-r}\le a_r\qquad(0\le r,\ 2r<D).
\end{equation}
The condition is vacuous for constants and for the zero polynomial,
whose natural degree we set to zero. The result holds even when a
coefficient sequence is not unimodal.

Write $\qi{m}=1+q+\cdots+q^{m-1}$ for $m\ge1$, and $\qi0=0$.
The four initial rows are $\delta_{0,k}$, and the recurrences are
\begin{align}
 S[n+1,k]&=S[n,k-1]+\qi{k}S[n,k],\notag\\
 c[n+1,k]&=c[n,k-1]+\qi{n}c[n,k],\label{eq:rec}\\
 S_B[n+1,k]&=S_B[n,k-1]+\qi{2k+1}S_B[n,k],\notag\\
 c_B[n+1,k]&=c_B[n,k-1]+\qi{2n+1}c_B[n,k].\notag
\end{align}
These are equations (2), (5), (8), and (9) of \cite{SS}'s first
arXiv version. Entries outside $0\le k\le n$ are zero; second-kind
recurrences need only be evaluated for $k\ge0$. Thus no negative
$q$-integer is needed. In particular, $S$ is the unshifted family,
not $q^{\binom{k}{2}}S[n,k]$.

Ma \cite{Ma} proves both bottom heaviness and monotonicity
on the upper half of the degree range. His public ChinaXiv deposit is dated
September 18, 2026. Principia Math in Progress
\cite[Lemmas 1--2 and Theorem 1]{PM} proves bottom heaviness for the
whole class of nonnegative combinations of products of $q$-integers.
The latter manuscript is dated September 12, 2026; its cited public
source was preserved on October 27, 2026. These dates distinguish a
manuscript date from a dated public copy. The proof below gives the
same closure statement through interval layers.

More precisely, Ma's Lemma 2 treats any nonnegative sum of symmetric
unimodal polynomials, Lemma 3 treats their products, and Lemma 4 applies
this mechanism to recurrences; his Proposition 1 treats elementary and
complete homogeneous symmetric functions on such alphabets
\cite[pp.~8--12]{Ma}. Thus the closure, upper-half monotonicity and
displayed complex-group consequences below also follow from his results.
Here the interval layers provide the rank injection and exact reflection
defect, including equality criteria and the symmetric-entry classification.

\section{Interval layers and an ambient degree}
For coefficients extended by zero outside $[0,d]$, symmetry and
unimodality mean $a_i=a_{d-i}$ and that the coefficients weakly
increase toward $d/2$. Under symmetry, unimodality is equivalent to
$a_i\le a_j$ whenever
$0\le i\le j$ and $i+j\le d$.

\begin{lemma}\label{lem:products}
Every finite product of $\qi{m}$ with $m\ge1$ is symmetric and
unimodal, with nonnegative coefficients and constant coefficient one.
\end{lemma}
\begin{proof}
A polynomial $f$ of degree $d$ with nonnegative symmetric unimodal
coefficients has the layer expansion
\begin{equation}\label{eq:layers}
 f(q)=\sum_{j=0}^{\lfloor d/2\rfloor}
 (a_j-a_{j-1})q^j\qi{d-2j+1},\qquad a_{-1}=0.
\end{equation}
Indeed, its coefficient at $i\in[0,d]$ is the telescoping sum
$\sum_{j=0}^{\min(i,d-i)}(a_j-a_{j-1})=a_{\min(i,d-i)}=a_i$.
All layer weights are nonnegative. For $a,b\ge1$, direct coefficient
counting gives
\begin{equation}\label{eq:interval}
 \qi a\qi b=\sum_{t=0}^{\min(a,b)-1}q^t\qi{a+b-1-2t}.
\end{equation}
Both sides have coefficient $\min(i+1,a,b,a+b-1-i)$ at
$0\le i\le a+b-2$. Each interval on the right is centered at
$(a+b-2)/2$, so the product is symmetric unimodal.
Multiplying each layer in \eqref{eq:layers} by $\qi m$ and applying
\eqref{eq:interval} produces symmetric unimodal layers with common
center $(d+m-1)/2$. Their sum is symmetric unimodal.
Induction starting from the empty product $1$ proves the assertion;
constant coefficients multiply to one.
\end{proof}

\begin{lemma}\label{lem:ambient}
Let $f$ have nonnegative, symmetric unimodal coefficients and degree
$d\le D$. Then $\coef{D-r}{f}\le\coef r f$ for every $0\le r$ with
$2r<D$.
\end{lemma}
\begin{proof}
If $D-r>d$, the left side is zero. Otherwise set $t=d-D+r$.
Then $0\le t\le r$ and $t+r=d-D+2r<d$. Symmetry followed by
increase toward the center gives
$\coef{D-r}{f}=\coef t f\le\coef r f$.
\end{proof}

\begin{theorem}\label{thm:closure}
Every polynomial in
\[
 \mathcal C=\left\{\sum_{\nu=1}^{N}\lambda_\nu
       \prod_{j=1}^{s_\nu}\qi{m_{\nu j}}:
       N,s_\nu\ge0,\ \lambda_\nu\in\mathbb R_{\ge0},\ m_{\nu j}\ge1\right\}
\]
is bottom heavy. If its actual degree is $D$, its coefficients also
satisfy $a_i\ge a_{i+1}$ for $\lceil D/2\rceil\le i<D$.
\end{theorem}
\begin{proof}
Discard zero weights. If nothing remains there is no strict-half
comparison. Otherwise nonnegativity and positive leading coefficients
give $D=\max_\nu d_\nu$, where $d_\nu=\sum_j(m_{\nu j}-1)$.
Apply Lemmas \ref{lem:products} and \ref{lem:ambient} to each product,
multiply by $\lambda_\nu$, and sum. This proves \eqref{eq:bh}.
For $i\ge\lceil D/2\rceil$, each summand's extended coefficient
sequence is on or beyond its own middle $d_\nu/2$, where it is
nonincreasing. Summing proves the last assertion.
\end{proof}

\begin{corollary}\label{cor:stirling}
All four unshifted families in \eqref{eq:rec} satisfy \eqref{eq:bh}
for every $n\ge0$ and $k\in\mathbb Z$.
The ordered second-kind families
$S^o[n,k]=\qi{k}!S[n,k]$ and
$S_B^o[n,k]=\qi{2k}!!S_B[n,k]$ have the same property for $k\ge0$.
\end{corollary}
\begin{proof}
The class $\mathcal C$ contains $0,1$ and is closed under addition
and multiplication by any $\qi m$, including $m=0$.
Induction on $n$ in all four recurrences puts every entry in
$\mathcal C$. The ordered multipliers are products of $q$-integers.
Apply Theorem \ref{thm:closure}.
\end{proof}

The same argument reproves Ma's symmetric-function consequence
\cite[Proposition 1]{Ma}: every elementary or complete homogeneous
symmetric polynomial evaluated at a finite alphabet
$\qi{m_1},\ldots,\qi{m_s}$ with positive integer lengths is bottom heavy,
because its monomial expansion is a nonnegative sum of products of these
entries. This includes the symmetric-function expressions in
\cite[Theorem 2.1]{SS}.

\begin{corollary}\label{cor:positive}
Every nonzero member of $\mathcal C$ has strictly positive coefficients
at all exponents from zero to its actual degree. In particular the
positivity clauses of \cite[Conjectures 7.4 and 7.5]{SS} hold for
nonzero entries.
\end{corollary}
\begin{proof}
A product of positive-length $q$-integers has positive coefficients
throughout its degree range: convolution of two positive sequences
supported on integer intervals has the same property. Choose a
positive-weight summand of maximum degree in a nonzero mixture.
Its coefficients are positive throughout the whole support of the sum.
\end{proof}

\section{Injections and equality}
\begin{corollary}\label{cor:injection}
For a mixture in $\mathcal C$ with integer weights, the product-of-chains
encoding admits an explicit injection from the objects of rank $D-r$
to those of rank $r$ whenever $2r<D$.
\end{corollary}
\begin{proof}
Label each summand and each of its copies separately. A product
$\prod_j\qi{m_j}$ counts the tuples $0\le x_j<m_j$ by rank $\sum_jx_j$.
Decompose this product of chains into symmetric saturated chains
as follows. For a rectangle $0\le x<a$, $0\le y<b$ with $a\le b$,
the chain indexed by $t=0,\ldots,a-1$ is
\[
 (t,0),(t,1),\ldots,(t,b-1-t),
 (t+1,b-1-t),\ldots,(a-1,b-1-t).
\]
These disjoint chains cover the rectangle, and their first and last
ranks sum to $a+b-2$. Interchange coordinates if $a>b$.
Applying this rectangle construction successively to each existing
chain and the next factor gives symmetric chains with rank intervals
$[j,d-j]$, where $d$ is the degree of the summand.

An element of rank $D-r$ lies in such a chain only if $D-r\le d-j$.
Since $d\le D$ and $r<D-r$, it follows that $j\le r\le d-j$.
Map the element to the unique element of rank $r$ on its chain,
retaining the summand and copy labels. Different chains or labels
cannot collide, and each chain has at most one element of rank $D-r$.
This proves injectivity. In particular the construction applies to
the labeled product expansions obtained from the four recurrences.
\end{proof}

\begin{corollary}\label{cor:defect}
Write $P=\sum_\nu\lambda_\nu f_\nu$, where the $f_\nu$ are nonzero
nonnegative symmetric unimodal polynomials, $\lambda_\nu>0$, and
$d_\nu=\deg f_\nu\le D=\deg P$.
Put $b_{\nu,j}=[q^j]f_\nu$, $b_{\nu,-1}=0$, and
$\delta_{\nu,j}=b_{\nu,j}-b_{\nu,j-1}$ for
$0\le j\le\lfloor d_\nu/2\rfloor$.
For every integer $0\le r<D/2$,
\begin{equation}\label{eq:defect}
 [q^r]P-[q^{D-r}]P=
 \sum_\nu\lambda_\nu
 \sum_{\substack{0\le j\le\lfloor d_\nu/2\rfloor\\
                  j\le r\le d_\nu-j\\ d_\nu-j<D-r}}
       \delta_{\nu,j}.
\end{equation}
Thus equality holds precisely when all the displayed $\delta_{\nu,j}$
vanish. Each positive displayed term gives a lower bound for the defect.
\end{corollary}
\begin{proof}
Use the layer expansion \eqref{eq:layers} in each summand.
A layer supported on $[j,d_\nu-j]$ that contains $D-r$ also contains
$r$, by the inequalities in the preceding proof. Its contribution to
the difference is therefore one exactly when it contains $r$ but not
$D-r$, and zero otherwise. All the layer weights are nonnegative,
which proves both the formula and the equality criterion.
\end{proof}

\begin{corollary}\label{cor:symmetric}
For $n\ge1$, among nonzero entries the symmetric cases are exactly
\[
\begin{array}{c|c}
 S[n,k]\text{ and }c[n,k]& k=1\text{ or }k=n\\
 S_B[n,k]\text{ and }c_B[n,k]& k=0\text{ or }k=n.
\end{array}
\]
In every other nonzero case the constant coefficient is strictly larger
than the leading coefficient. The initial row and all zero entries are symmetric.
\end{corollary}
\begin{proof}
Use $S[n,k]=h_{n-k}(\qi1,\ldots,\qi k)$ and
$c[n,k]=e_{n-k}(\qi1,\ldots,\qi{n-1})$, and the type $B$
formulas in \cite[Theorem 2.1]{SS}.
Every product in these expansions has constant and leading coefficient one.
Consequently the constant coefficient of the sum counts all products,
whereas the leading coefficient counts just those of maximum degree.
For $S[n,k]$ with $2\le k<n$, the products $\qi1^{n-k}$ and
$\qi k^{n-k}$ have different degrees; for $S_B[n,k]$ the same argument
uses $\qi1$ and $\qi{2k+1}$ when $1\le k<n$.
For either first-kind family, the elementary products of a fixed
size strictly between zero and the alphabet size have different degrees,
because the alphabet lengths are strictly increasing.
In the remaining cases the polynomial is constant or a single product
of $q$-integers, hence symmetric by Lemma~\ref{lem:products}.
\end{proof}

\section{Other source consequences}
For $m\ge2$ write $[km]!_m=\prod_{i=1}^k\qi{im}$.
The following statements concern the unshifted polynomials defined by
Sagan and Swanson in \cite{Complex}.
\begin{corollary}\label{cor:complex}
For $0\le k\le n$, the following polynomials are bottom heavy,
positive throughout their nonempty supports, and nonincreasing on
the upper half of their degree ranges:
\begin{align*}
 S[m,n,k]&=h_{n-k}(\qi1,\qi{m+1},\ldots,\qi{km+1}),\\
 (-1)^{n-k}s[m,n,k]&=e_{n-k}(\qi1,\qi{m+1},\ldots,\qi{(n-1)m+1}),\\
 \widetilde S[m,n,k]&=h_{n-k}(\qi{m-1},\qi{2m-1},\ldots,\qi{(k+1)m-1}),
\end{align*}
as well as their products with $[km]!_m$.
The Chan--Rhoades polynomial $S^o_{CR}[m,n,k]$ has these properties;
the Hilbert polynomial whose coefficient reversal it equals is top heavy.
\end{corollary}
\begin{proof}
The formulas are \cite[Theorem 3.2(16), Section 6.1,
equation (20), and Proposition 5.10]{Complex}.
The last identifies $S^o_{CR}=[km]!_m S[m,n,k]$ as a reversed Hilbert
polynomial. All the displayed alphabets have positive integer lengths,
so their symmetric-function expansions and the factorial multipliers
belong to $\mathcal C$. Apply Theorem~\ref{thm:closure} and
Corollary~\ref{cor:positive}. Reversing coefficients changes the
bottom-heavy inequalities into the top-heavy inequalities
$a_r\le a_{D-r}$ for $2r<D$.
\end{proof}

These deductions do not cover the subtractive formula (17) of
\cite{Complex}, arbitrary signed first-kind coefficients, or every
irreducible complex reflection group. They concern the displayed
families, not the conjectural super-coinvariant interpretations.

\begin{corollary}\label{cor:alternating}
The alternating sum in Sagan--Swanson Conjecture 7.9 satisfies, for $n\ge0$,
\begin{equation}\label{eq:alternating}
 \sum_{k=0}^n(-1)^{n-k}q^{2(n-k)}\qi{k}!S[n,k]-1
   =(1-q)(1+q)^n-1+q^{n+1}.
\end{equation}
For $n\ge1$ its coefficient at $q^j$, $1\le j\le n$, is
$\binom nj-\binom n{j-1}$; coefficients at zero and $n+1$ vanish.
The nonzero coefficients are positive below $(n+1)/2$, negative above
it, and paired with their negatives by $j\mapsto n+1-j$.
For $n\ge2$, factoring out the initial $q$ makes the absolute-coefficient
polynomial palindromic in the ordinary degree convention.
\end{corollary}
\begin{proof}
Use the classical falling-factorial identity
\[
 x^n=\sum_{k=0}^n S[n,k]\prod_{j=0}^{k-1}(x-\qi j),
\]
as in \cite[equation (3.12)]{DZ}. Substituting $x=-q^{-1}$
and $x=-q^{-2}(1+q)$ in the rational function field gives respectively
\[
 \sum_k(-1)^{n-k}q^{n-k}\qi{k}!S[n,k]=1,\qquad
 \sum_k(-1)^{n-k}q^{2(n-k)}\qi{k+1}!S[n,k]=(1+q)^n.
\]
Let $F$ be the sum on the left of \eqref{eq:alternating} before
subtracting one. Since $(1-q)\qi{k+1}=1-q^{k+1}$, these two
identities imply $F-q^{n+1}=(1-q)(1+q)^n$.
This is an identity of polynomials, hence includes $q=0,1$.
The coefficient formula follows by the binomial theorem. Finally
$\binom nj>\binom n{j-1}$ exactly when $j<(n+1)/2$;
binomial symmetry gives the asserted sign pairing. For $n=0$ the
polynomial is zero.
\end{proof}

\section{Scope and mathematical coverage}
The symmetry convention in Conjecture 7.9 needs care. If $P_n$ denotes
the left side of \eqref{eq:alternating}, then $P_2=q-q^2$.
Its absolute-coefficient polynomial $q+q^2$ is not palindromic
with respect to its actual degree two. The exact symmetry is
$q^{n+1}P_n(1/q)=-P_n(q)$ for $n\ge1$; after factoring out $q$,
the absolute coefficients are palindromic. Thus the sign assertions
hold, and support normalization makes the symmetry assertion precise.

The degree bound and exponent-zero convention are essential.
For $(1+q)^2=1+2q+q^2$, comparison at the false degree bound $D=1$
would require $2\le1$.
Also $S[2,2]=1$, whereas the shifted polynomial $q$ fails
\eqref{eq:bh} at $r=0$. Neither example contradicts the theorem.
For the unshifted families, $S[4,3]=3+2q+q^2$ illustrates bottom
heaviness without symmetry. The type $B$ examples $S_B[6,4]$ and
$c_B[7,5]$ in \cite[Section 7.5]{SS} are bottom heavy but not unimodal.

\begin{center}\small
\begin{tabularx}{\textwidth}{@{}lX@{}}
\toprule
Source item in \cite{SS}, arXiv v1 & Coverage \\
\midrule
Conjecture 7.6 & Reproved; prior proofs: Ma (2026) and Principia Math in Progress (2026). \\
Ordered families & Reproved; consequences of Ma's product and sum lemmas. \\
Theorem 2.1 & Symmetric-function consequence reproved; prior proof: Ma (2026). \\
Conjectures 7.4, 7.5 & Positivity clauses proved for nonzero entries; coefficient and parity log-concavity not covered. \\
Conjecture 7.9 & Sign assertions proved; symmetry corrected to $P_n/q$. The actual-degree palindrome reading fails at $n=2$, where $P_2=q-q^2$. \\
Conjectures 6.5, 6.7 & Prior proofs cited: Rhoades--Wilson for the Hilbert series, then Angarone et al. for the type $A$ basis \cite{RW,ACKMR}. \\
Conjectures 6.15, 6.17 & Type $B$ super-coinvariant Hilbert series and bases not covered. \\
Sections 7.2--7.4 & Major-index/basis analogues, solving the $q$-difference equation, and a sign-reversing proof of Theorem 2.2 not covered. \\
Section 7.6, Conjecture 7.8 & Asymptotic normality and limit shapes not covered. \\
Section 7.1; \cite{Complex} & Displayed complex-group consequences of Ma's Proposition 1 reproved in Corollary~\ref{cor:complex}; other cases not covered. \\
\cite{PM}, Tasks 2--4 & Product-of-chains injection proved; defect and symmetric-case formulas proved; displayed complex-group extensions proved. Other encodings and analogues not covered. \\
\bottomrule
\end{tabularx}
\end{center}

\paragraph{Mathematical coverage.}
The Lean declaration
\begin{quote}\small\ttfamily
D5.S3.Combinatorics.QStirling.\\
SaganSwansonBottomHeavy.result
\end{quote}
in the
\href{https://github.com/the-omega-institute/trureturing/blob/13f4288c2940bd2aa4020cec789956b9265769f5/D5/S3/Combinatorics/QStirling/SaganSwansonBottomHeavy.lean}{pinned \texttt{SaganSwansonBottomHeavy.lean} source},
proves precisely the four-family natural-coefficient assertion of
Corollary \ref{cor:stirling}, including all integer $k$ and the
natural-degree convention for zero. The real-weight closure,
upper-half monotonicity, injections, equality formulas, ordered and
complex-group consequences, and alternating sum have the written proofs above.
The accompanying exact-integer checker tests all four recurrences
for $0\le n\le50$, $-3\le k\le53$, symmetric-function identities
for $0\le n\le8$, and finite product and unequal-degree mixture
examples. It additionally checks the chain injections and equality
formulas, the displayed complex-group recurrences for $2\le m\le5$
and $n\le10$, and \eqref{eq:alternating} for $n\le30$.
These finite checks illustrate the universal proofs.

\begin{thebibliography}{9}
\bibitem{SS}
B. E. Sagan and J. P. Swanson,
\emph{$q$-Stirling numbers in type $B$},
European Journal of Combinatorics \textbf{118} (2024), 103899.
\href{https://doi.org/10.1016/j.ejc.2023.103899}{doi:10.1016/j.ejc.2023.103899}.
Original posing version:
\href{https://arxiv.org/abs/2205.14078v1}{arXiv:2205.14078v1} (2022),
Section 7.5, Conjecture 7.6, p.~40.
\bibitem{Ma}
S.-M. Ma,
\emph{Log-concavity and log-convexity of Stirling columns, and
bottom-heaviness of $q$-Stirling numbers},
ChinaXiv:202609.00184v1, deposited September 18, 2026.
\href{https://chinaxiv.org/abs/202609.00184v1}{Version record},
\href{https://doi.org/10.12074/202609.00184}{doi:10.12074/202609.00184}.
\href{https://www.researchgate.net/publication/414355049_Log-concavity_and_log-convexity_of_Stirling_columns_and_bottom-heaviness_of_q-Stirling_numbers}{Author-uploaded full text},
Lemmas 2--4, Theorems 4--5 and Proposition 1, pp.~8--12.
\bibitem{PM}
Principia Math in Progress,
\emph{Bottom heaviness of the four $q$-Stirling families:
a solution to MathDB problem 354788},
note dated September 12, 2026, Lemmas 1--2 and Theorem 1.
\href{https://github.com/PrincipiaMathOrg/math-solutions/blob/47371706c400c79d06e3fb62190fcf6f7a725c85/site-backup/mathdb-354788/mathdb-354788-solution.tex}{Version-specific scientific source}.
\bibitem{Complex}
B. E. Sagan and J. P. Swanson,
\emph{Stirling numbers for complex reflection groups},
Annals of Combinatorics (2025).
\href{https://doi.org/10.1007/s00026-025-00751-4}{doi:10.1007/s00026-025-00751-4}.
Formulas cited from
\href{https://arxiv.org/abs/2408.13874v1}{arXiv:2408.13874v1} (2024).
\bibitem{DZ}
M.-J. Ding and J. Zeng,
\emph{Some identities involving $q$-Stirling numbers of the second kind in type $B$},
Electronic Journal of Combinatorics \textbf{31}(1) (2024), P1.36.
\href{https://doi.org/10.37236/12147}{doi:10.37236/12147}.
\bibitem{ACKMR}
R. Angarone, P. Commins, T. Karn, S. Murai and B. Rhoades,
\emph{Superspace coinvariants and hyperplane arrangements},
\href{https://arxiv.org/abs/2404.17919v2}{arXiv:2404.17919v2} (2026),
Corollary 7.2; first version 2024.
\bibitem{RW}
B. Rhoades and A. Wilson,
\emph{The Hilbert series of the superspace coinvariant ring},
Forum of Mathematics, Pi \textbf{12} (2024), e16, Theorem 5.3.
\href{https://doi.org/10.1017/fmp.2024.14}{doi:10.1017/fmp.2024.14}.
\end{thebibliography}
\medskip\noindent\small
\copyright\ 2026 Haobo Ma and Wenlin Zhang.
Licensed under Creative Commons Attribution 4.0 International
(CC BY 4.0).
\end{document}
